% main_lrcbart.tex : JSM 2026 deck rewritten from the LRC-BART revision
% (revision/05-writing/main_draft.tex, 2026-09-09). The earlier deck describing
% MAP-BART is kept as main.tex. Numbers: revision/05-writing/tables/quoted_numbers.txt,
% 04-application/application_report.md. Figures: inserts/lrc_fig_sc4_map.pdf,
% inserts/lrc_fig_ess_map.pdf (copies of revision/05-writing/figures/).
\documentclass[11pt,aspectratio=169]{beamer}

\usepackage[english]{babel}
\usepackage[T1]{fontenc}
\usepackage{amsmath,amssymb,bm,mathtools}
\usepackage{booktabs,multirow}
\usepackage{graphicx}
\usepackage[round]{natbib}
\usepackage{xcolor}
\usepackage{tikz}
\usetikzlibrary{arrows.meta,positioning}

\newcommand{\Normal}{\mathrm{N}}
\newcommand{\E}{\mathbb{E}}
\newcommand{\Prb}{\mathbb{P}}
\newcommand{\bx}{\bm{x}}
\newcommand{\given}{\,\vert\,}
\newcommand{\target}{N_{\mathrm{target}}}
\newcommand{\Ber}{\mathrm{Bernoulli}}

\mode<presentation>{
  \usetheme[progressbar=frametitle,sectionpage=none]{metropolis}

  \definecolor{BayesGreen}{RGB}{0,80,0}
  \definecolor{BayesLight}{RGB}{232,242,232}
  \definecolor{BayesGold}{RGB}{176,125,0}
  \setbeamercolor{structure}{fg=BayesGreen}
  \setbeamercolor{alerted text}{fg=red!80!black}
  \setbeamercolor{example text}{fg=BayesGreen}
  \setbeamercolor{palette primary}{fg=white,bg=BayesGreen}
  \setbeamercolor{palette secondary}{fg=BayesGreen,bg=BayesLight}
  \setbeamercolor{palette tertiary}{fg=white,bg=BayesGreen!80}
  \setbeamercolor{block title}{bg=BayesGreen!22,fg=black}
  \setbeamercolor{block body}{bg=BayesGreen!5,fg=black}
  \setbeamertemplate{blocks}[rounded][shadow=false]
  \setbeamerfont{block body}{size=\footnotesize}
  \setbeamerfont{block body alerted}{size=\footnotesize}
  \setbeamerfont{block title}{size=\small}
  \setbeamerfont{block title alerted}{size=\small}
  \setbeamersize{text margin left=0.55cm,text margin right=0.55cm}
  \setbeamertemplate{navigation symbols}{}
  \setbeamerfont{title}{size=\Large}
  \setbeamerfont{subtitle}{size=\small}
  \setbeamertemplate{footline}{
    \leavevmode
    \hbox{%
      \begin{beamercolorbox}[wd=.34\paperwidth,ht=2.25ex,dp=1ex,leftskip=1em]{author in head/foot}
        \scriptsize JSM 2026 \enspace\textbullet\enspace Yuan Ji
      \end{beamercolorbox}%
      \begin{beamercolorbox}[wd=.66\paperwidth,ht=2.25ex,dp=1ex,rightskip=1em]{title in head/foot}
        \raggedleft\scriptsize \insertshorttitle\hspace*{1.5em}\insertframenumber
      \end{beamercolorbox}%
    }
    \vskip0pt
  }
}

\title[LRC-BART]{Locally Robust Commensurate BART\\for Borrowing External Controls}
\subtitle{Local compatibility and borrowing calibration}
\author{Yuan Ji}
\institute{University of Chicago\\Department of Public Health Sciences}
\date{Joint Statistical Meetings 2026}

\begin{document}

\begin{frame}[plain]
  \vspace{0.7em}
  {\usebeamerfont{title}\usebeamercolor[fg]{title}\inserttitle\par}
  \vspace{0.5em}
  {\usebeamerfont{subtitle}\usebeamercolor[fg]{subtitle}\insertsubtitle\par}
  \vspace{1em}
  {\color{red!70}\rule{\textwidth}{0.4pt}}
  \vfill
  {\usebeamerfont{author}\insertauthor\par}
  \vspace{0.45em}
  {\small Collaborators: Yunxuan Zhang (UChicago), Gu Mi (Sanofi), Ji Lin (Sanofi)\par}
  \vspace{0.35em}
  {\usebeamerfont{date}\insertdate\par}
  \vspace{0.9em}
  {\usebeamerfont{institute}\insertinstitute\par}
  \vspace{0.7em}
\end{frame}

\begin{frame}{Overview}
  \begin{enumerate}
    \item Established borrowing priors use one study-level parameter. Compatibility with an external control can vary across patient profiles.
    \item LRC-BART adds a sparse discrepancy ensemble to the pooled BART surface, with a spike-and-slab commensurate prior on each discrepancy leaf.
    \item We calibrate a variance-based prior ESS for the standardized control mean. Its ceiling depends on the RWD fit at trial profiles.
    \item Simulations with regional shifts, and a 30-patient single-arm myeloma trial with a registry control.
  \end{enumerate}
\end{frame}

\section{Motivation}
\begin{frame}{External controls and transportability}
  \begin{columns}[T,onlytextwidth]
    \column{0.48\textwidth}
      \begin{block}{Potential gain}
        \begin{itemize}
          \item More control observations
          \item Greater precision
          \item Smaller concurrent control arms
        \end{itemize}
      \end{block}
    \column{0.48\textwidth}
      \begin{alertblock}{Why pooling can mislead}
        \begin{itemize}
          \item Different eligibility and care pathways
          \item Different covariate distributions
          \item Measured and unmeasured confounding
        \end{itemize}
      \end{alertblock}
  \end{columns}
  \vspace{0.5em}
  \centering
  \textbf{External controls may agree with the trial in one patient group and differ in another. A single borrowing parameter cannot distinguish these groups.}
\end{frame}

\begin{frame}{Established priors use a study-level borrowing parameter}
  \small
  \begin{columns}[T,onlytextwidth]
    \column{0.48\textwidth}
      \begin{block}{Meta-analytic predictive prior \citep{Neuenschwander2010MAP,schmidli2014robust}}
        Study parameters are exchangeable, $\theta_j\mid\mu,\tau^2\sim\Normal(\mu,\tau^2)$; the robust version adds a vague component that suppresses borrowing under conflict.
      \end{block}
      \begin{block}{Commensurate prior \citep{Hobbs2011,Hobbs2012}}
        $\theta\mid\theta_0\sim\Normal(\theta_0,1/\tau)$ with a spike-and-slab hyperprior on the commensurability $\tau$.
      \end{block}
    \column{0.48\textwidth}
      \begin{block}{Power prior \citep{IbrahimChen2000}}
        Historical likelihood raised to $a_0$.
      \end{block}
      \begin{alertblock}{Common feature}
        One heterogeneity or discount scale applies to the whole population, without depending on the patient's covariates. With a single historical study the three forms nearly coincide \citep{RoverFriede2025}.
      \end{alertblock}
  \end{columns}
  \vspace{0.4em}
  \centering
  \textbf{We allow borrowing to vary across covariate regions learned from the data.}
\end{frame}

\section{Model}
\begin{frame}{The trial-control mean adds a discrepancy to the external mean}
  \small
  RCT controls ($D_i=1$) and RWD controls ($D_i=2$), covariates $\bx_i$, outcome $y_i$:
  \[
    \boxed{\;y_i=f(\bx_i)+\mathbb{I}\{D_i=1\}\,g(\bx_i)+\varepsilon_i,\qquad \varepsilon_i\sim\Normal(0,\sigma^2_{D_i})\;}
  \]
  \begin{columns}[T,onlytextwidth]
    \column{0.48\textwidth}
      \begin{block}{$f$: external-control response mean}
        Standard BART, $H_f=50$ trees, fitted to both sources. The RWD follow $f$.
      \end{block}
    \column{0.48\textwidth}
      \begin{block}{$g$: discrepancy ensemble}
        Small and shallow, $H_g=5$ trees, split prior $(0.5,3)$, informed by the RCT controls only. The RCT controls follow $f_1=f+g$.
      \end{block}
  \end{columns}
  \vspace{0.4em}
  \begin{itemize}\footnotesize
    \item The decomposition follows Bayesian causal forests \citep{Hahn2020BCF}, with source membership replacing treatment.
    \item The discrepancy is $g=f_1-f$. Shrinking it toward zero centers the RCT-control mean $f_1$ at the external-control mean $f$.
  \end{itemize}
\end{frame}

\begin{frame}{Each discrepancy leaf has a spike-and-slab prior}
  \footnotesize
  The zero-centered mixture describes each leaf contribution to $g=f_1-f$:\par\vspace{-0.5em}
  $\theta_{h\ell}$ is the contribution of leaf $(h,\ell)$ to $g$:
  \[
    \theta_{h\ell}\given w,\tau_0^2,\tau_1^2
    \stackrel{\mathrm{iid}}{\sim} w\,\Normal(0,\tau_0^2)+(1-w)\,\Normal(0,\tau_1^2).
  \]
  Equivalently, $z_{h\ell}\given w\stackrel{\mathrm{iid}}{\sim}\Ber(w)$ selects a zero-mean normal with variance $\tau_{1-z_{h\ell}}^2$. We use $w\sim\mathrm{Beta}(1,1)$, $\tau_0^2\sim\text{scaled-Inv-}\chi^2(\nu_0,s_0^2)$ truncated to $(0,\tau_1^2)$, and $\tau_1^2$ fixed by the range rule.
  \begin{columns}[T,onlytextwidth]
    \column{0.48\textwidth}
      \begin{block}{Spike, $z=1$: the commensurate component}
        Shrinks a discrepancy contribution toward zero, drawing the RCT response mean toward $f(\bx)$. The ESS calibration sets the leaf scale $\tau_0$.
      \end{block}
    \column{0.48\textwidth}
      \begin{block}{Slab, $z=0$: the vague component}
        A free local offset on the outcome scale, corresponding to a robust-MAP vague component within a leaf.
      \end{block}
  \end{columns}
  \vspace{0.3em}
  \begin{block}{Special cases}
    $w=1$: commensurate BART (C-BART), one learned scale. \quad $w=1,\ \tau_0^2\to0$: BART on the pooled controls (BART-CP). \quad $w=0$: a free source offset, the tree analogue of BART with the source indicator (BART-PP).
  \end{block}
\end{frame}

\begin{frame}{The induced prior is centered at $f(\bx)$}
  \small
  \begin{itemize}
    \item At a fixed profile $\bx$, each discrepancy tree contributes one leaf:
      $g(\bx)=\sum_{h=1}^{H_g}\theta_{h,\ell_h(\bx)}$, where $\ell_h(\bx)$ is the reached leaf.
    \item Integrating the leaf indicators gives a mixture over the number $m$ of reached spike leaves:
  \end{itemize}
  \[
  f_1(\bx)\given f,T_g,w,\tau_0^2,\tau_1^2
  \sim\sum_{m=0}^{H_g}\binom{H_g}{m}w^m(1-w)^{H_g-m}
  \Normal\!\left(f(\bx),m\tau_0^2+(H_g-m)\tau_1^2\right).
  \]
  \begin{columns}[T,onlytextwidth]
    \column{0.48\textwidth}
      \begin{block}{All spike: $m=H_g$}
        Variance $H_g\tau_0^2$: the commensurate component shrinks $f_1(\bx)$ toward $f(\bx)$.
      \end{block}
    \column{0.48\textwidth}
      \begin{block}{All slab: $m=0$}
        Variance $H_g\tau_1^2$: a larger departure is permitted, with the same center $f(\bx)$.
      \end{block}
  \end{columns}
  \vspace{0.3em}
  \begin{itemize}\footnotesize
    \item $T_g$ denotes the discrepancy partitions; every component is centered at the external mean $f(\bx)$.
    \item This is a prior for the latent control mean. Observation variance is separate.
  \end{itemize}
\end{frame}

% Diagram: one illustrative discrepancy tree, not a fitted partition.
\begin{frame}{The discrepancy changes across patient profiles}
\centering
\begin{tikzpicture}[
  x=1cm,y=1cm,
  split/.style={draw,rounded corners=1pt,fill=white,minimum size=7pt,inner sep=1pt},
  spike/.style={draw=UChicagoMaroon,rounded corners=2pt,fill=UChicagoMaroon!8,inner sep=5pt,font=\small},
  slab/.style={draw=AccentBlue,rounded corners=2pt,fill=AccentBlue!8,inner sep=5pt,font=\small},
  edge label/.style={font=\scriptsize,fill=white,inner sep=2pt},
  note/.style={font=\footnotesize,align=left},
  callout/.style={draw=UChicagoMaroon,dashed,thick,rounded corners=3pt,fill=white,inner sep=7pt,align=left,font=\footnotesize}]
\path[use as bounding box] (0,0) rectangle (14.5,5.65);
\node[anchor=west,font=\small] at (0.15,5.3)
  {\textbf{RWD controls:} $f(\bx)$ \qquad \textbf{RCT controls:} $f(\bx)+g(\bx)$};
\node[anchor=west,note] at (0.15,4.68) {One tree in $g$ \; (illustrative splits)};
\node[split] (root) at (3.4,4.15) {};
\node[spike] (left) at (1.45,2.95) {$\theta_{h1}$};
\node[split] (right) at (5.3,2.95) {};
\node[slab] (middle) at (4.25,1.72) {$\theta_{h2}$};
\node[spike] (last) at (6.45,1.72) {$\theta_{h3}$};
\draw (root) -- (left) node[midway,edge label,above left] {$x_1<c_1$};
\draw (root) -- (right) node[midway,edge label,above right] {$x_1\ge c_1$};
\draw (right) -- (middle) node[midway,edge label,left] {$x_2<c_2$};
\draw (right) -- (last) node[midway,edge label,right] {$x_2\ge c_2$};
\node[callout,anchor=north west,text width=6.45cm] (sum) at (0.15,0.85)
  {Each profile reaches one leaf per tree.\\
   Add those contributions: $g(\bx)=\sum_{h=1}^{H_g}\theta_{h,\ell_h(\bx)}$.};
\draw[AccentBlue,dashed,thick,-{Stealth}] (sum.north) -- (middle.south);
\node[note,anchor=north west,text width=6.05cm] at (8.05,4.7)
  {\textbf{$f$: the external response mean}\\
   Fitted jointly to both control sources.};
\node[spike,anchor=north west,text width=5.8cm,align=left,font=\footnotesize] at (8.05,3.70)
  {\textbf{Spike: $z_{h\ell}=1$}\\
   $\theta_{h\ell}\sim\Normal(0,\tau_0^2)$\\
   Keeps the RCT mean close to $f(\bx)$.};
\node[slab,anchor=north west,text width=5.8cm,align=left,font=\footnotesize] at (8.05,2.02)
  {\textbf{Slab: $z_{h\ell}=0$}\\
   $\theta_{h\ell}\sim\Normal(0,\tau_1^2)$,\quad $\tau_1^2>\tau_0^2$\\
   Allows a local departure from $f(\bx)$.};
\node[note,anchor=north west,text width=6cm,font=\scriptsize] at (8.05,0.45)
  {$z_{h\ell}\sim\Ber(w)$. Leaf indicators vary. The parameters $w,\tau_0^2,\tau_1^2$ are shared.};
\end{tikzpicture}
\end{frame}


\begin{frame}{Separate ensembles allow different shrinkage priors}
  \small
  \begin{itemize}
    \item Separate ensembles let us choose different partition and shrinkage priors for the external mean $f$ and the discrepancy $g$.
    \item We use $H_f=50$ trees for the outcome surface and $H_g=5$ shallow trees for the discrepancy.
    \item At fixed leaf variances and partitions, fewer discrepancy trees reduce the total all-spike variance at a profile. The slab range rule also depends on $H_g$.
  \end{itemize}
  \begin{block}{What the prior-density comparison establishes}
    Appendix G.6 minimizes over integer numbers of slab leaves and retains contributions from the remaining spike leaves. It compares joint prior densities at constrained modes. Posterior indicator probabilities also depend on the likelihood and partitions.
  \end{block}
  \footnotesize The simulations assess this specification. They do not establish that shared-leaf models cannot detect regional discrepancies.
\end{frame}

\begin{frame}{Tree moves use exact leaf marginals}
  \footnotesize
  Conditionally on $(\tau_0^2,\tau_1^2,w)$ the leaf prior is a finite Gaussian mixture, which gives closed-form factors for each tree move:
  \[
    M_f=(1+\lambda_f^2A)^{-1/2}\exp\Bigl\{\frac{\lambda_f^2B^2}{2(1+\lambda_f^2A)}\Bigr\},
    \qquad
    M_g=\frac{w\,\phi(\bar y_1;0,\tau_0^2+v_1)+(1-w)\,\phi(\bar y_1;0,\tau_1^2+v_1)}{\phi(\bar y_1;0,v_1)},
  \]
  with $A,B$ the precision-weighted counts and sums of both sources in an $f$ leaf, and $\bar y_1$, $v_1=\sigma_1^2/n_1$ the RCT residual mean and its variance in a $g$ leaf.
  \begin{itemize}
    \item Grow, prune and change moves use the exact collapsed ratios, so the sampler targets the exact posterior without a surrogate likelihood.
    \item Leaf draws are conjugate; $z$ is drawn first, then $\theta$; $\tau_0^2$ and $w$ update from the leaf means and states.
    \item For right-censored outcomes, log-normal AFT data augmentation leaves the other updates unchanged.
    \item One sweep costs a standard BART on $n_1+n_2$ observations plus a five-tree ensemble on $n_1$.
  \end{itemize}
\end{frame}

\section{Borrowing calibration}
\begin{frame}{Prior effective sample size at the estimand}
  \footnotesize
  For $\mu=N^{-1}\sum_i f_1(\bx_i)$, $V_\mu^f$ is the RWD-only posterior variance. The working calibration is:
  \[
    \mathrm{ESS}_0(s_0^2)=\hat\sigma_1^2\,\E_{\tau_0^2}\Bigl[\frac{1}{V_\mu^f+c_gH_g\tau_0^2}\Bigr].
  \]
  \begin{itemize}
    \item $c_g$ averages squared leaf proportions before inversion; the scale law here is untruncated.
    \item A normal prior gives the ELIR variance ratio \citep{neuenschwander2020}; this BART/scale-mixture summary is not exact marginal ELIR.
    \item Ceiling: $\hat\sigma_1^2/V_\mu^f$, including BART regularization. Equality to scaled $n_2$ requires one fixed tree, matched proportions and a flat prior.
    \item The Pr-rule requires the target in 95\% of blocks. A target at or above the whole-chain ceiling instead triggers a smallest-scale override.
  \end{itemize}
  \begin{block}{Realized ESS diagnostic}
    With target leaf proportions $a_{h\ell}$, use $V_g=\sum_{h,\ell}a_{h\ell}^2\{z_{h\ell}\tau_0^2+(1-z_{h\ell})\tau_1^2\}$, then average $\sigma_1^2/(V_\mu^f+V_g)$. Standardization can span several leaves per tree.
  \end{block}
\end{frame}

\begin{frame}{Conditional theoretical guarantees}
  \footnotesize
  Leaf results condition on $f$, the other discrepancy trees, partitions, scales and $0<w<1$. Write $v=\sigma_1^2/n$ for a leaf's RCT residual-mean variance.
  \begin{block}{Theorem 1: bounded local influence}
    Relative to the slab estimate, $\sup_{\bar y}|I(\bar y)|\le\sqrt{2v\max\{\log C,\tfrac12\}}$, and the shift vanishes at a Gaussian rate under growing conflict. The C-BART and BART-CP leaf shifts are linear and unbounded.
  \end{block}
  \begin{block}{Theorem 2: local detection and map consistency}
    The leaf spike probability decreases with $|\bar y|$. With fixed partitions, $f$ fixed, a representable discrepancy and growing local counts, $\mathcal B_c(x)\to\mathbb I\{|d|<c\}$ for $|d|\ne c$. A fixed positive spike scale need not identify compatibility through its indicator.
  \end{block}
  \begin{block}{Theorem 3: the working calibration}
    Under the untruncated law, $\mathrm{ESS}_0$ is continuous, convex and strictly decreasing from its ceiling to zero. Truncation gives a positive lower limit. Grid-selection consistency requires the stated moment and separation conditions; ceiling overrides are separate.
  \end{block}
\end{frame}

\begin{frame}{The borrowing map}
  \footnotesize
  At each patient profile $x$, we report
  \[
    \mathcal B_c(x)=P\bigl(|g(x)|<c\given\text{data}\bigr),
  \]
  the posterior probability that the local discrepancy is below a threshold $c$ fixed in advance on the outcome scale, together with the posterior mean and interval of $g(x)$ and the pointwise ESS.
  \begin{itemize}
    \item $c=0.5$ in the simulations, one third of the RCT residual SD for Gaussian outcomes.
    \item The all-spike probability $B(x)$ is kept as a diagnostic: its prior scale is $w^{H_g}$ and at a fixed positive spike scale, it need not identify compatibility as RCT counts grow.
    \item In single-arm mode the data do not inform $g$, so the compatibility map is flat. The local ESS map describes variation across profiles.
  \end{itemize}
  \begin{alertblock}{Conditions for the map result}
    The map results require fixed $g$ partitions, a representable discrepancy, growing local counts and $f$ fixed; the result does not cover the posterior over partitions. The simulations show separation of shifted and compatible regions under the full posterior.
  \end{alertblock}
\end{frame}

\section{Simulation evidence}
\begin{frame}{Simulation design}
  \footnotesize
  \begin{columns}[T,onlytextwidth]
    \column{0.40\textwidth}
      \textbf{Common design}
      \begin{itemize}\setlength\itemsep{1pt}
        \item $R=500$ paired replicates
        \item $p=10$ covariates, nonlinear surface
        \item $300$ RWD controls
        \item Gaussian: RCT $200{:}100$, $\sigma_1=1.5$, $\sigma_2=1.8$
        \item Survival: RCT $225{:}75$, log-normal AFT, about a quarter censored
        \item Targets: absolute $50,75,100$; fractional $0.25,0.5,0.9$ of the ceiling
      \end{itemize}
    \column{0.58\textwidth}
      \textbf{Five scenarios}
      \begin{itemize}\setlength\itemsep{1pt}
        \item \textbf{Sc1} no confounding: source at random
        \item \textbf{Sc2} measured: source depends on $(X_5,X_6)$
        \item \textbf{Sc3} unmeasured: source and outcome depend on $\bm U$, observed as proxies at $\rho\in\{-0.5,0,0.5\}$
        \item \textbf{Sc4} regional shift: RWD surface raised by $\delta_2\in\{0.5,1,2\}$ outside $\mathcal R=\{X_5>2\}$, half the population; secondary version on the null covariate $X_7$
        \item \textbf{Sc5} global shift: raised by $\delta_2$ everywhere; borrowing should be suppressed
      \end{itemize}
      \vspace{0.2em}
      {\scriptsize Illustrative leaf threshold: $d_{1/2}\approx0.5$. Comparators: regression/BART pooling variants (PP includes a source indicator), PSCL, MAP, hierarchies, C-BART and LRC-BART ($w=0$). CAHB \citep{Jin2023CAHB}: transferred analysis, first $200$ paired Gaussian replicates; variants in the supplement.}
  \end{columns}
\end{frame}

\begin{frame}{Gaussian outcomes: ATE bias, RMSE, coverage}
  \centering
  \scriptsize
  \setlength{\tabcolsep}{3.2pt}
  \renewcommand{\arraystretch}{0.9}
  \resizebox{\textwidth}{!}{%
  \begin{tabular}{lrrr|rrr|rrr|rrr|rrr}
    \toprule
    & \multicolumn{3}{c|}{Sc1} & \multicolumn{3}{c|}{Sc3, $\rho=0$} & \multicolumn{3}{c|}{Sc4, $\delta_2=1$} & \multicolumn{3}{c|}{Sc4, $\delta_2=2$} & \multicolumn{3}{c}{Sc5, $\delta_2=2$} \\
    Method & Bias & RMSE & Cov & Bias & RMSE & Cov & Bias & RMSE & Cov & Bias & RMSE & Cov & Bias & RMSE & Cov \\
    \midrule
    LM-PP   & $-$0.00 & 0.34 & 0.94 & 0.04 & 0.35 & 0.95 & $-$0.01 & 0.34 & 0.95 & $-$0.02 & 0.34 & 0.95 & $-$0.04 & 0.34 & 0.94 \\
    BART-NP & 0.00 & 0.22 & 0.91 & $-$0.00 & 0.29 & 0.94 & 0.00 & 0.22 & 0.91 & 0.00 & 0.22 & 0.91 & 0.00 & 0.22 & 0.91 \\
    BART-CP & 0.01 & 0.14 & 0.96 & 0.98 & 1.00 & 0.00 & $-$0.34 & 0.37 & 0.35 & $-$0.68 & 0.70 & 0.00 & $-$1.39 & 1.40 & 0.00 \\
    BART-PP & 0.01 & 0.21 & 0.94 & 0.03 & 0.27 & 0.94 & $-$0.00 & 0.21 & 0.95 & $-$0.01 & 0.21 & 0.96 & $-$0.01 & 0.21 & 0.95 \\
    C-BART  & 0.01 & 0.15 & 0.95 & 0.23 & 0.38 & 0.87 & $-$0.21 & 0.27 & 0.73 & $-$0.23 & 0.39 & 0.66 & $-$0.32 & 0.65 & 0.76 \\
    \textbf{LRC-BART (100)} & \textbf{0.01} & \textbf{0.17} & \textbf{0.94} & \textbf{0.04} & \textbf{0.28} & \textbf{0.94} & $\bm{-0.04}$ & \textbf{0.22} & \textbf{0.89} & \textbf{0.00} & \textbf{0.20} & \textbf{0.92} & $\bm{-0.00}$ & \textbf{0.20} & \textbf{0.92} \\
    \bottomrule
  \end{tabular}}
  \vspace{0.3em}
  \begin{block}{Agreement, global conflict and regional shifts}
    Agreement (Sc1): LRC-BART borrows, RMSE $0.172$ against $0.205$ for BART-PP; the pooling models have lower RMSE here. Global disagreement (Sc3, Sc5): it suppresses borrowing, with bias about $0.04$ and $0.00$, where C-BART and BART-CP have substantial bias. Regional shift: near-zero empirical bias at $\delta_2=2$; at $\delta_2=1$ it borrows across the boundary, bias $-0.04$, coverage $0.89$.
  \end{block}
\end{frame}

\begin{frame}{The borrowing map under a regional shift}
  \centering
  \includegraphics[width=0.74\textwidth]{inserts/lrc_fig_sc4_map.pdf}
  \vspace{-0.3em}
  \begin{block}{LRC-BART at target 100, first 40 replicates}
    At $\delta_2=2$ the map is $0.02$ outside $\mathcal R$ and $0.7$ inside, and $g$ is $-1.8$ and $-0.15$ against $-2$ and $0$; the transition spans the unit around the boundary, with about four RCT controls per $0.1$ of $X_5$. At $\delta_2=1$ the estimated discrepancy captures part of the shift ($-0.6$ / $-0.25$). Under Sc5 the map is $0.01$ everywhere.
  \end{block}
\end{frame}

\begin{frame}{Regional conflict remains difficult}
  \footnotesize
  \begin{block}{Agreement and global conflict}
    Sc1: RMSE $0.172$ versus $0.205$ for BART-PP; realized ESS $53$, prior ESS $82$, ceiling $127$. Under global conflict, realized ESS falls to $1.6$ (Sc3) and $0.2$ (Sc5). Sc3 bias is $0.012$ above BART-PP (paired SE $0.001$); the leaf theorem does not establish the cause of this difference.
  \end{block}
  \begin{alertblock}{Surface error under regional shifts}
    At shift $2$ ($1.3$ residual SDs), compatible-region surface RMSE over all RCT profiles is $0.824$ versus $0.785$. On control profiles away from the boundary, it is $0.747$ versus $0.723$ (difference $0.024$, SE $0.005$, $R=200$). At shift $1$ (two thirds of a residual SD), bias is $-0.04$, coverage $0.89$, and RMSE $0.223$ versus $0.207$.
  \end{alertblock}
  \begin{block}{Regional ESS depends on the standardized mean}
    At shift $2$, realized regional ESS is $9.5$ in the compatible region and $0.44$ outside ($R=200$). Each uses its own standardized mean; they are not additive. The compatibility map reports a posterior probability.
  \end{block}
  \footnotesize The CAHB comparison assesses our transfer of kernel-local estimation; it does not establish a general method ranking. Sc2's ceiling falls from $127$ to $23$, so absolute targets are capped.
\end{frame}

\begin{frame}{Survival outcomes: efficiency under agreement}
  \centering
  \footnotesize
  RMST ratio at $t^\ast=3$; $R=500$ simulated trials
  \vspace{0.1em}

  \setlength{\tabcolsep}{6pt}
  \renewcommand{\arraystretch}{0.9}
  \begin{tabular}{lrrr|rrr}
    \toprule
    & \multicolumn{3}{c|}{LRC-BART (100)} & \multicolumn{3}{c}{BART-PP} \\
    Scenario & Bias & RMSE & Cov & Bias & RMSE & Cov \\
    \midrule
    Sc1: agreement & 0.003 & 0.097 & 0.94 & $-$0.023 & 0.106 & 0.94 \\
    Sc2: measured confounding & 0.018 & 0.100 & 0.93 & $-$0.002 & 0.097 & 0.95 \\
    Sc3: unmeasured confounding & 0.031 & 0.146 & 0.94 & $-$0.005 & 0.137 & 0.92 \\
    Sc4: regional shift & $-$0.037 & 0.108 & 0.90 & $-$0.059 & 0.113 & 0.91 \\
    Sc5: global shift & 0.001 & 0.112 & 0.93 & $-$0.041 & 0.110 & 0.93 \\
    \bottomrule
  \end{tabular}
  \vspace{0.2em}

  {\scriptsize BART-PP: pooled trial and external controls, with a source indicator as a predictor.}
  \begin{block}{Efficiency and robustness}
    Under agreement, LRC-BART improves RMST precision and power ($0.28$ versus $0.17$ for BART-PP). RMST RMSE remains close to BART-PP across the studied conflict scenarios.
  \end{block}
  {\scriptsize Median-ratio estimation is less favorable under unmeasured confounding; regional RMST coverage is $0.90$. Sc3: $\rho=0$; Sc4--5: $\delta_2=1$. Full results: paper and Web Table S2.}
\end{frame}

\begin{frame}[t]{Single-arm trials cannot learn a control discrepancy}
  \footnotesize
  Trial contributes treated subjects only ($n_T=200$ Gaussian; $30$ or $200$ survival); the $300$ RWD controls are the sole control source; LRC-BART in single-arm mode with $w=1$.
  \begin{columns}[T,onlytextwidth]
    \column{0.49\textwidth}
      \begin{block}{No confounding}
        BART-CP and LRC-BART have Gaussian RMSE $0.16$; the survival RMST ratio at $n_T=30$ has RMSE $0.16$ to $0.19$ with coverage $0.92$ to $0.97$.
      \end{block}
    \column{0.49\textwidth}
      \begin{alertblock}{Measured confounding}
        Every method is biased: Gaussian ATE bias $1.68$ (LM-CP), $0.41$ (BART-CP), $0.42$ (LRC-BART); there are no internal controls to identify a discrepancy, and RWD overlap is limited.
      \end{alertblock}
  \end{columns}
  \vspace{0.3em}
  \begin{itemize}
    \item The single-arm ceiling is $27$ to $34$ under Sc2, so an absolute target of $100$ is capped and LRC-BART approximates pooling with $w=1$; the fractional target still widens the interval (coverage $0.82$ to $0.99$).
    \item With $w<1$, slab variance $(1-w)H_g\tau_1^2$ adds uncertainty: SD $1.4$ for the ATE and $10$ to $34$ for the median ratio. Zero-mean $g$ preserves linear log-time means, but survival-ratio means also change. Primary analysis: $w=1$.
    \item Hierarchical regressions: ATE SD $8$ under the stated hyperprior.
  \end{itemize}
\end{frame}

\section{EloKRd application}
\begin{frame}{EloKRd and UCMM data}
  \begin{columns}[T,onlytextwidth]
    \column{0.48\textwidth}
      \begin{block}{EloKRd phase II trial}
        \begin{itemize}
          \item $n=30$, all received EloKRd
          \item 13 PFS events; 10 deaths
          \item 50\% high-risk cytogenetics; 30\% ASCT
        \end{itemize}
      \end{block}
    \column{0.48\textwidth}
      \begin{block}{UCMM registry control}
        \begin{itemize}
          \item $n=253$ after cohort restrictions
          \item Standard induction regimens
          \item Five baseline covariates plus ASCT after induction
        \end{itemize}
      \end{block}
  \end{columns}
  \vspace{0.3em}
  \begin{alertblock}{Harmonization and the scope of the comparison}
    Inconsistent non-Hispanic indicator coding changes extrapolation. Harmonization raises the working ESS ceilings from $12.6$/$14.6$ to $63.7$/$36.7$ (PFS/OS). ASCT is recorded after induction, so this adjustment supports a descriptive comparison, not a total causal treatment effect.
  \end{alertblock}
  \centering\footnotesize Estimand: standardized 5-year RMST ratio (neither trial median is reached). Kaplan--Meier: $1.015$ for PFS, $0.90$ for OS.
\end{frame}

\begin{frame}{EloKRd results: 5-year RMST ratio, posterior mean and 95\% interval}
  \centering
  \scriptsize
  \setlength{\tabcolsep}{6pt}
  \begin{tabular}{lcc}
    \toprule
    Harmonized coding & PFS & OS \\
    \midrule
    Kaplan--Meier, unadjusted & 1.01 (0.84, 1.22) & 0.90 (0.79, 1.03) \\
    Complete-pooling AFT & 0.87 (0.69, 1.06) & 0.80 (0.65, 0.94) \\
    Hierarchical AFT & 1.23 (0.74, 1.90) & 1.21 (0.77, 1.70) \\
    AFT-BART on UCMM, no discrepancy & 0.97 (0.79, 1.16) & 0.90 (0.78, 1.02) \\
    \midrule
    \textbf{LRC-BART, $\target=30$} & \textbf{0.97 (0.77, 1.23)} & \textbf{0.90 (0.77, 1.04)} \\
    LRC-BART, $\target=15$ & 0.99 (0.75, 1.40) & 0.91 (0.76, 1.11) \\
    LRC-BART, $0.9\times$ ceiling & 0.97 (0.79, 1.15) & 0.90 (0.77, 1.04) \\
    \midrule
    \textit{Earlier analysis, original coding:} MAP-AFT-BART, 30 & 1.08 (0.66, 16.4) & 0.97 (0.74, 4.77) \\
    \bottomrule
  \end{tabular}
  \vspace{0.3em}
  \begin{block}{Descriptive comparison conditional on recorded profiles}
    At target $30$, $P(\text{ratio}>1)$ is $0.37$ (PFS) and $0.08$ (OS); both nonparametric intervals include one. Lower targets widen intervals. The fitted PFS distribution has a $99.9$th percentile of $1.6$ over $8{,}000$ draws, but the earlier sampler is unavailable, so we cannot isolate the cause of the tail difference.
  \end{block}
\end{frame}

\begin{frame}{The local ESS map at the 30 trial profiles}
  \centering
  \includegraphics[height=0.56\textheight]{inserts/lrc_fig_ess_map.pdf}
  \vspace{-0.2em}
  \begin{block}{Registry support at the trial profiles}
    $\mathrm{ESS}(x)=\hat\sigma_1^2/\{V_f(x)+H_g\bar\tau_0^2\}$ ranges from $5$ to $17$ for PFS and $6$ to $15$ for OS: highest for White/ASCT/standard-risk profiles; lowest for the one Hispanic patient, Other/Unknown race and high-risk profiles without ASCT. $\bar\tau_0^2$ is the fitted mean scale. Target $30$ concerns the standardized log-time mean; local values are not additive.
  \end{block}
\end{frame}

\section{Takeaways}
\begin{frame}{Main findings}
  \begin{enumerate}
    \item The RCT control mean is centered on the external mean through a spike-and-slab discrepancy ensemble.
    \item Closed-form leaf marginals support posterior-invariant tree updates; this does not establish finite-chain mixing.
    \item Variance-based ESS calibration concerns the standardized control mean. The calculation specifies its ceiling and approximation convention.
    \item Conditional leaf influence and map results require the stated fixed-surface and partition assumptions.
    \item Gaussian Sc1 RMSE is about $16\%$ lower than BART-PP's. Regional conflict produces informative compatibility summaries, but also surface losses and excessive borrowing at moderate shifts.
  \end{enumerate}
\end{frame}

\begin{frame}{Limitations and next steps}
  \begin{itemize}
    \item Full-posterior map consistency over learned partitions is unproved; the Sc4 surface loss persists away from the boundary.
    \item Calibration uses a working precision summary. Exact marginal ELIR would require a different calculation and new calibration analyses.
    \item Multiple external sources require joint calibration. Additive ceilings have not been established.
    \item A single-arm study cannot learn the control discrepancy. The application remains a descriptive comparison with post-induction ASCT adjustment.
    \item Confirmatory use requires operating-characteristic evaluation, including Type I error and power, and a pre-specified role for the map.
  \end{itemize}
  \begin{block}{How to interpret the borrowing summaries}
    The compatibility map, prior pointwise ESS and realized regional ESS answer different questions. These summaries do not give an observed count of external patients used at each profile.
  \end{block}
\end{frame}

\begin{frame}[allowframebreaks]{References}
  \fontsize{6.5pt}{7.4pt}\selectfont
  \bibliographystyle{unsrtnat}
  \bibliography{ref_lrcbart}
\end{frame}

\end{document}
